\documentclass[border=10pt]{standalone}
\usepackage{tikz}
\usetikzlibrary{arrows.meta, positioning, shapes.geometric, calc}

\begin{document}

\definecolor{blockcolor}{RGB}{0, 120, 200}
\definecolor{arrowcolor}{RGB}{80, 80, 80}
\definecolor{fbcolor}{RGB}{220, 50, 30}
\definecolor{sumcolor}{RGB}{40, 160, 40}
\begin{tikzpicture}[
    block/.style={
      rectangle, draw=blockcolor, line width=1.5pt,
      fill=blockcolor!8, rounded corners=3pt,
      minimum width=26mm, minimum height=14mm,
      font=\bfseries, text=blockcolor, align=center
    },
    % Summing-junction circle diameter is ~40% of block height (Block
    % Diagram Style convention), not 71% as before (10mm / 14mm block).
    % 加え合わせ点の円の直径はブロック高さの約4割（ブロック線図の流儀）。
    % 以前は 10mm / ブロック高さ14mm = 71% だった。
    sum/.style={
      circle, draw=sumcolor, line width=1.5pt,
      fill=sumcolor!8, minimum size=6mm,
      font=\normalsize\bfseries, text=sumcolor
    },
    arrow/.style={-{Stealth[length=3mm]}, line width=1.2pt, arrowcolor},
    lbl/.style={font=\small, above, text=arrowcolor},
  ]

  % Blocks
  \node[block] (target) at (0, 0) {Target\\Rate};
  \node[sum]   (sum)    at (3, 0) {$\Sigma$};
  \node[block] (kp)     at (6, 0) {P Control\\$K_p \times e$};
  \node[block] (mixer)  at (9.5, 0) {Motor\\Mixer};
  \node[block] (drone)  at (13, 0) {StampFly\\Drone};

  % Forward path
  \draw[arrow] (target) -- node[lbl] {$r$} (sum);
  \draw[arrow] (sum) -- node[lbl] {$e$} (kp);
  \draw[arrow] (kp) -- node[lbl] {$u$} (mixer);
  \draw[arrow] (mixer) -- (drone);

  % Output line with feedback tap point
  \draw[arrow] (drone.east) -- ++(2.5, 0) node[lbl, right] {response};
  \coordinate (fb_tap) at ($(drone.east)+(1.2, 0)$);
  \fill[arrowcolor] (fb_tap) circle (2.5pt);

  % Feedback path (branches from output line, not from inside drone)
  \node[block, fill=fbcolor!8, draw=fbcolor, text=fbcolor]
    (gyro) at (8, -3) {Gyro\\Sensor};
  \draw[arrow, fbcolor] (fb_tap) -- (fb_tap |- gyro) -- (gyro);
  \draw[arrow, fbcolor] (gyro) -- (3, -3) -- (sum);

  % Sum signs: "+" above the left-entering arrow, "-" to the RIGHT of the
  % bottom-entering feedback arrow (Block Diagram Style convention).
  % 加え合わせ点の符号: 「+」は左から入る矢印の上、「-」は下から入る
  % フィードバック矢印の右側に置く（ブロック線図の流儀）。
  % Fix: the "-" was at 210 deg (left of the upward arrow) with anchor
  % "north east" (glyph extending further left) -- it sat on the WRONG
  % (left) side of the arrow. Moved to 330 deg (right of the arrow) with
  % anchor "north west" (glyph extending right of the point).
  % 修正: 以前「-」は210度（上向き矢印の左側）・anchor=north east
  % （グリフがさらに左に伸びる）で、矢印の誤った（左）側にあった。
  % 330度（矢印の右側）・anchor=north west（グリフが右に伸びる）へ移動。
  \node[font=\scriptsize, sumcolor, anchor=south east] at ($(sum.north)+(150:0.3)$) {$+$};
  \node[font=\scriptsize, fbcolor, anchor=north west] at ($(sum.south)+(330:0.3)$) {$-$};

  % Labels
  \node[font=\small, arrowcolor, above] at ($(drone.east)+(1.2, 0.5)$) {angular rate};
  \node[font=\small, fbcolor, below] at (8, -3.8) {feedback};

\end{tikzpicture}
\end{document}
